\chapter{Introduction}
\label{chap:introduction}

\section{Background}
\label{section:background}

In the current software development landscape, the industry has seen a surge in tools designed to accelerate the development lifecycle.
However, a significant disparity remains between the tools available for low-level code implementation and those focused on the high-level detection and explanation of architectural integrity.

Currently, the primary solutions for architectural analysis fall into two categories: Large Language Models (LLMs) and specialized static analysis platforms such as SonarQube.
While these tools provide baseline insights, they present several limitations as effective detection systems for modern development teams:

\begin{itemize}
    \item \textbf{Operational Cost and Accessibility:} Enterprise-grade analysis tools often incur high licensing costs,
    making them inaccessible for smaller teams or individual researchers. Furthermore, cloud-dependent LLMs require continuous financial investment
    and a stable internet connection for consistent analysis.

    \item \textbf{Metric-Centric Detection:} Existing tools are highly effective at generating quantitative metrics (e.g., cyclomatic complexity)
    but often fail to detect \textit{structural} patterns. They typically report on localized code quality rather than identifying systemic architectural smells
    that emerge from the interactions between components.

    \item \textbf{The Interpretive Burden:} Most detection tools shift the entire responsibility of interpretation back to the developer.
    Without clear, natural language guidance, a developer must manually bridge the gap between a flagged numeric metric and the underlying architectural flaw it represents.

    \item \textbf{LLM Hallucinations and Bias:} Generic LLM-based solutions are prone to hallucinatory outputs when analyzing complex codebases.
    This can lead to the detection of non-existent issues or inappropriate suggestions that do not suit the specific structural constraints of the system.
\end{itemize}

Ultimately, traditional tools remain focused on "low-level" detection, such as syntax styling and local logic errors.
This leaves a critical opening for a system like \textbf{\usevar{\srsNickname}} that graph-based architectural detection with local execution and high explainability.

\section{Problem Statement}
\label{section:problem-statement}

The modern systems are increasingly complex and are often developed at a rapid pace to meet market demands.
This acceleration frequently comes at the cost of architectural integrity, leading to several critical challenges that \textbf{\usevar{\srsNickname}} aims to address:

\subsection{Neglect of High-Level Architectural Analysis}
\label{subsection:architectural-neglect}
While modern development workflows are saturated with automated testing and continuous integration,
there is a notable lack of tooling dedicated specifically to the \textbf{architectural aspect} of software.
Existing static analysis tools often focus on low-level code quality—such as naming conventions,
code style, and local design patterns—leaving high-level architectural analysis to manual inspection or fragmented, ad-hoc tools.
Consequently, structural issues often remain undetected until they manifest as significant technical debt.

\subsection{Language Fragmentation and Tooling Silos}
\label{subsection:language-fragmentation}
The shift toward polyglot development and multi-language repositories has created a ``tooling silo'' effect.
Most robust analysis solutions are language-specific, which creates significant friction when attempting to enforce architectural consistency across a diverse codebase.
There is a distinct lack of cross-language frameworks capable of providing a unified architectural perspective regardless of the underlying programming languages used.

\subsection{The Explainability Gap}
\label{subsection:explainability-gap}
Even when advanced analysis is performed, a significant gap often exists between the complexity of the detection and the clarity of the resulting explanation.
Existing tools frequently flag structural issues without providing the necessary context or natural language justification. For a \textbf{software architect},
a simple ``warning'' is insufficient; without a metric-backed explanation of \textit{why} a particular design is problematic, it is difficult to reason about or prioritize improvements.
While developers may understand the source code, they may lack the specialized knowledge to interpret raw metric values.
\textbf{\usevar{\srsNickname}} aims to bridge this gap by providing clear,
natural language explanations supported by quantitative metrics.

\section{Solution Overview}
\label{section:solution-overview}

\textbf{\usevar{\srsNickname}} aims to provide clear, natural language explanations for architectural smell detection.
Using \textbf{Tree-sitter} to help parse and understand the structure of source code into a standardized \textbf{TOML-based Intermediate Representation (IR)},
which is then used to detect architectural smells using \textbf{GraphSAGE}.
The result of classification will then be used to generate natural language explanations for the detected smells backed by quantitative metrics.

The \textbf{TOML-based IR} approach will ensure decoupling of analysis model and source code, allowing for easy adaptation to new languages and technologies.

\subsection{Features}
\label{subsection:features}

\begin{description}[style=multiline, leftmargin=3.5cm, font=\normalfont]
    \item[\textbf{Multi-Language Parsing}] Automated extraction of structural data from source code into a standardized \textbf{TOML-based Intermediate Representation (IR)}.
    \item[\textbf{Graph-Based Modeling}] Transformation of software components and their dependencies into a directed graph format optimized for \textbf{Graph Neural Network (GNN)} analysis.
    \item[\textbf{Architectural Smell Detection}] Intelligent identification of complex structural \textit{anti-patterns}, including Cyclic Dependencies, God Components, and Hub-like dependencies.
    \item[\textbf{Semantic XAI Explanation}] Generation of natural language justifications using \textbf{Explainable AI (XAI)} to describe why specific structural patterns were flagged as architectural smells.
    \item[\textbf{Interactive Visualization}] A dynamic graphical interface for exploring the system's architecture, allowing users to visually trace dependencies and identified smells.
\end{description}

\section{Target User}
\label{section:target-user}

<TIP:
The target user in a software project refers to the specific
group or demographic of individuals for whom the software is designed and
developed. Identifying the target user is a crucial step in the software
development process as it helps the development team tailor the software to
meet the needs, preferences, and requirements of that particular user group.
Understanding the characteristics, behaviors, and expectations of the target
users is essential for creating a user-friendly and effective software solution.

Here are some key aspects related to defining the target user in a software project:

Demographics: This includes factors such as age, gender,
occupation, education level, and other demographic characteristics. Different
age groups or professional backgrounds may have distinct preferences and
requirements when it comes to software usability.

Skill Level: Consideration of the users' technical proficiency and
familiarity with similar software or technology. The level of technical expertise
can influence the complexity of the user interface, the need for tutorials or
documentation, and other user support features.

Industry or Domain: For software solutions designed for specific
industries or domains, understanding the unique challenges, workflows, and
terminology within that industry is crucial. Tailoring the software to meet
industry-specific needs is often necessary.
/>

\section{Terminology}
\label{section:terminology}

<TIP:
Terminology refers to the specific language, jargon, or
specialized vocabulary used to describe concepts, ideas, or subjects within a
particular field or domain. The use of terminology is often essential for clarity
and precision, especially in books that cover technical, scientific, academic, or
specialized topics.
/>
